% --- ejercicio_1.tex ---
\hypertarget{ej1}{}
\hypertarget{ej1}{}
\noindent \textbf{Ejercicio 1.} Determinar los valores de verdad de las siguientes proposiciones cuando el valor de verdad de $a$, $b$ y $c$ es verdadero y el de $x$ e $y$ es falso.

\begin{multicols}{2}
\begin{enumerate}[label=\alph*)]
    \item $(\neg x \vee b)$
    \item $((c \vee (y \wedge a)) \vee b)$
    \item $\neg(c \vee y)$
    \item $\neg(y \vee c)$
    \item $(\neg(c \vee y) \leftrightarrow (\neg c \wedge \neg y))$
    \item $((c \vee y) \wedge (a \vee b))$
    \item $(((c \vee y) \wedge (a \vee b)) \leftrightarrow (c \vee (y \wedge a) \vee b))$
    \item $(\neg c \wedge \neg y)$
\end{enumerate}
\end{multicols}

\vspace{0.2cm}
{\color{gray}\hrule}
\vspace{0.4cm}

\textbf{Resolución:}

\textbf{a)} $(\neg x \vee b)$

Para dar una demostración del valor de verdad de esta proposición, primero identificamos 
los valores dados por el enunciado. Sabemos que el valor de $b$ es verdadero ($V$) y el de $x$ es falso ($F$). 

Como ya tenemos sus valores fijos, basta con armar una pequeña tabla de valores paso a paso. Primero calculamos la negación de $x$ ($\neg x$) y luego resolvemos la disyunción lógica ($\vee$) con $b$:

\begin{center}
\begin{tabular}{|c|c|c|c|}

$x$ & $\neg x$ & $b$ & $(\neg x \vee b)$ \\
F & V & V & \textbf{V} \\
\end{tabular}
\end{center}

Y así comprobamos que la proposición $(\neg x \vee b)$ es \textbf{Verdadera}.

\vspace{0.5cm}
\aportar
\hrule
\vspace{0.5cm}
\aportar
\hrule
\vspace{0.5cm}
